\documentclass[11pt,a4paper]{article}
\usepackage[margin=2.6cm]{geometry}
\usepackage[T1]{fontenc}
\usepackage{lmodern}
\usepackage[letterspace=180]{microtype}
\usepackage{booktabs}
\usepackage{amsmath,amssymb}
\usepackage{parskip}
\usepackage{xcolor}
\usepackage{titlesec}
\usepackage{caption}
\usepackage{graphicx}
\usepackage[hidelinks]{hyperref}

\definecolor{bcnavy}{RGB}{24,38,68}
\definecolor{bcgold}{RGB}{198,160,74}
\definecolor{bcteal}{RGB}{23,110,120}
\definecolor{bcgrey}{RGB}{95,95,95}

\titleformat{\section}{\bfseries\large\color{bcnavy}}{\thesection.}{0.6em}{}
\titleformat{\subsection}{\bfseries\color{bcteal}}{}{0em}{}
\captionsetup{font={footnotesize},labelformat=empty,textfont={color=bcgrey}}

\begin{document}

% ---------------------------------------------------------------- cover
\begin{titlepage}
\centering
\vspace*{3.2cm}
\includegraphics[width=5cm]{bc_logo.pdf}

\vspace{1.6cm}
{\color{bcnavy}\rule{7cm}{0.6pt}}

\vspace{0.9cm}
{\LARGE\bfseries\color{bcnavy} Algorithmic Trading\\[0.15cm]
Model Development\par}
\vspace{0.35cm}
{\large\color{bcnavy} A summary of the models designed, built and deployed:\\
two statistical pricing engines, a derivative-free optimisation framework,\\
and the validation programme that governs deployment\par}

\vspace{0.9cm}
{\color{bcgold}\rule{7cm}{1.0pt}}

\vspace{1.5cm}
{\bfseries Bayesian Capital}\\[0.2cm]
{\color{bcgrey} Systematic Trading Research}\\[0.9cm]
{\itshape\color{bcgrey} Internal document --- strictly confidential}\\[0.3cm]
{\color{bcgrey} 6 October 2026}
\vfill
\end{titlepage}

% ---------------------------------------------------------------- header
\begin{center}
{\Large\bfseries\color{bcnavy} Algorithmic Trading Model Development}\\[0.15cm]
{\color{bcnavy} a summary of the models behind the nine-stock book}\\[0.25cm]
{\small\color{bcgrey} Bayesian Capital --- Systematic Trading Research
\;$\bullet$\; current book: TSM $\cdot$ VRT $\cdot$ VST $\cdot$ AVGO $\cdot$ MU
$\cdot$ GM $\cdot$ VLO $\cdot$ CF $\cdot$ MRVL}
\end{center}
\vspace{0.1cm}
{\color{bcnavy}\hrule height 0.8pt}
\vspace{0.4cm}

\section{What this note covers}

Over the course of 2026 we designed, built, tested and deployed a systematic
trading system. It runs a live nine-stock equity book out of a
spreadsheet-and-broker production environment, and every entry price it
places is computed by one of two statistical models built in-house. This note
summarises the mathematics of that development for readers who do not need to
work through it in detail; the full treatment, with the derivations, the
calibration evidence and the out-of-sample validation programme, is held in
the firm's internal research documentation.

The development falls into four distinct bodies of work: the two pricing
engines, a derivative-free global optimisation framework for fitting them,
and the validation programme that decides what is deployed and what is not.
Each is summarised in turn.

\section{The strategy and its capital mechanics}

The book harvests short-term volatility. It buys weakness at a
model-computed discount and sells a fixed premium into the bounce that
usually follows. Each trading day, for each stock, the engines produce a
single object: a limit bid $B_t$, computed causally from data up to the
previous close. Everything downstream of that bid --- the fill logic, the
premium exit, the stop-loss, commissions and interest on idle cash --- is
shared between the two engines, so they can be compared directly and blended
freely.

Capital goes to work in staggered tranches. For $n$ tranches of $K/n$ each,
with $M_j$ the $j$-th tranche's terminal return multiple over its round
trips,
\begin{equation}
\Pi=\sum_{j=1}^{n}\frac{K}{n}\bigl(M_j-1\bigr)=K\bigl(\bar M-1\bigr),
\end{equation}
so the book's profit is governed by the \emph{average} multiple across
start-dates. This identity is what lets the two engines and nine names share
one pooled capital loop, and it is the reason capital allocation was treated
as a research question in its own right (\S6).

\section{Engine I: a Bayesian local-linear-trend Kalman filter}

The first engine treats a stock's fair value as something to be inferred
rather than guessed at. It carries a two-dimensional state, a level $m_t$ and
a trend $b_t$, observed through noisy closes:
\begin{equation}
m_t=m_{t-1}+b_{t-1}+\eta^{(m)}_t,\qquad
b_t=b_{t-1}+\eta^{(b)}_t,\qquad
C_t=m_t+\varepsilon_t .
\end{equation}
What separates this from a textbook filter is that all three noise variances
are scaled to the day's trading range $R_t=\bar H_t-\bar L_t$:
\begin{equation}
\operatorname{Var}\bigl(\eta^{(m)}_t\bigr)=(\varphi_L R_t)^2,\qquad
\operatorname{Var}\bigl(\eta^{(b)}_t\bigr)=(\psi R_t)^2,\qquad
\operatorname{Var}(\varepsilon_t)=(\lambda R_t)^2,
\end{equation}
so the filter automatically trusts the tape less on turbulent days. The
update is the standard Kalman recursion in precision form: the posterior
mean is a precision-weighted average of the prior and the new datum, with
the gain $K_t=V_t^-/(V_t^-+r_t)$ computed afresh every day rather than set
by hand. The entry bid sits a fitted number of \emph{predictive} standard
deviations below the one-step forecast, floored at the day's open and at a
cap below the all-time high $P_t$:
\begin{equation}
\sigma_t=\sqrt{V_t+q_t+r_t},\qquad
B_t^{\text{LLT}}=\min\bigl((\hat m_{t-1}+\hat b_{t-1})-k\,\sigma_{t-1},\;O_t,\;(1-c)P_t\bigr).
\end{equation}
Because $\sigma_t$ widens when the filter is less certain, the discount
deepens in turbulence of its own accord. There are no hand-set weights
anywhere in the engine.

\section{Engine II: an Ornstein--Uhlenbeck mean-reversion sleeve}

The second engine was built to be different on purpose. Where the Kalman
filter projects the trend forward, this one fades deviations from a slow
long-run mean, so its mistakes land on different days. It models price as
the mean-reverting diffusion
\begin{equation}
dX_t=\theta(\mu-X_t)\,dt+\sigma\,dW_t,
\end{equation}
which, sampled daily, is an exact AR(1) about the mean,
$X_t-\mu=\phi(X_{t-1}-\mu)+\varepsilon_t$ with $\phi=e^{-\theta\Delta t}$.
The persistence $\phi$ is estimated each day as the least-squares AR(1)
coefficient over a rolling window of $W$ closes and clamped to $[0,0.99]$.
The clamp is substantive, not cosmetic, since values outside it assert
either phantom anti-persistence or a non-stationary random walk. The
discount scale is the standard deviation of the fitted AR(1)
\emph{residuals},
\begin{equation}
e_s = d_s-\hat\phi_t\,d_{s-1},\qquad
s_t=\sqrt{\tfrac{1}{W-1}\textstyle\sum_s (e_s-\bar e_t)^2},
\end{equation}
rather than of the price levels, and that distinction turned out to matter.
A level-based scale conflates trend with innovation; identifying the error,
correcting it and evidencing the correction was part of this year's work,
and it remains the only change to a live signal that has survived our
out-of-sample tests. The bid takes the same form as Engine~I's: a $k$-sigma
discount below the one-step forecast
$f_t=\bar C_t+\hat\phi_t(C_{t-1}-\bar C_t)$, floored by the open and the
peak cap.

The sleeve's role in the book was also re-measured during the year. The two
engines' returns correlate at $0.57$ unconditionally and at $0.90$ or higher
when both hold stock, so the sleeve is carried as a timing diversifier and a
source of return, not as the near-independent hedge an earlier analysis had
suggested. Establishing that, including showing that a previously reported
near-zero correlation was the artefact of an idle sleeve, is part of the
documented work.

\section{Parameter estimation and optimisation}

Each of the nine names carries a ten-dimensional parameter vector,
\begin{equation}
\theta=(\underbrace{\lambda,\varphi_L,\psi,k,\pi,c}_{\text{Kalman engine}},\;
\underbrace{W,b_{\text{OU}},\pi_{\text{OU}},c_{\text{OU}}}_{\text{OU engine}})\in\Theta\subset\mathbb R^{10},
\end{equation}
fitted by maximising a backtest objective subject to a minimum trade count.
The objective surface is awkward in a specific way: it is piecewise
constant. A fill is the step indicator $\mathbf 1\{L_t\le B_t(\theta)\}$, so
the gradient is zero almost everywhere and jumps at the fill boundaries.
Gradient and quasi-Newton methods are useless here, and the fit instead uses
differential evolution, a population-based global optimiser that needs no
derivatives. The search is seeded by a Sobol sequence with the incumbent
parameters injected as one member of the population, and it runs under a
fixed random seed so that every fit is reproducible.

To avoid fragile knife-edge optima, the objective is averaged over a
neighbourhood of each candidate: with multiplicative perturbations
$\rho^{(p)}\sim\mathcal U(0.97,1.03)$,
\begin{equation}
\tilde S(\theta)=\tfrac12 S(\theta)+\tfrac12\frac1P\sum_{p=1}^{P}
S\bigl(\Pi_\Theta[\theta\odot\rho^{(p)}]\bigr),
\end{equation}
which penalises spikes and favours plateaus. The same thinking carries
through to reporting: planning figures are quoted at the median of the
parameter neighbourhood, not at the fitted point.

\section{The validation programme}

The largest single body of work is neither of the models. It is the testing
regime that decides what gets deployed.

Whether parameters should be re-fitted to recent data was tested directly,
across in-sample, fit-one-year-test-the-next and rolling walk-forward
designs. Re-fitted parameters beat the incumbents in 6 of 40 out-of-sample
folds on the daily model and in 1 of 18 on a weekly variant that has since
been retired. That result now governs our change control: parameters are
frozen, and only specification corrections, as opposed to re-fits, are
eligible for deployment.

Allocation was researched with the same discipline. Equal weighting across
names was tested against inverse-volatility, Sharpe-proportional and
return-proportional schemes and beat them all by roughly 10 points out of
sample. The capital split between the two engines was tested across its full
range, on return and on risk-adjusted objectives, with the buffer re-fitted
inside each fold; that analysis is what moved the book from a 75:25 split to
the flat 50:50 it runs today.

The backtests themselves are not simulations of a formula alone. A Python
replica of the production engine reproduces the live workbook's cached
results exactly, and same-day round trips are verified against 1-minute or
5-minute intraday bars before they are counted.

Finally, the full book, including its deployed protection gate (a
breadth-conditional 200-day moving-average rule), was replayed through the
2022 bear market on intraday data, to measure drawdown behaviour in a regime
the fitting history does not contain.

\section{Deployment and scale}

The system runs live on nine US-listed names, each on both engines with its
own fitted ten-parameter vector, pooled in a single capital loop and
operated daily through a spreadsheet-and-broker production stack with
automated data refresh. The modelled activity level is on the order of five
hundred fills a year. The development is documented across a suite of
internal papers --- a full methodology paper, an operations manual, a
book-construction and ranking memorandum, a bear-market protection study and
a strategy evolution record --- together with the research codebase of
optimisers, replay engines, validation harnesses and operational tooling,
from which every figure quoted here can be reproduced.

\vspace{0.4cm}
{\color{bcnavy}\hrule height 0.8pt}
\vspace{0.2cm}
{\footnotesize\color{bcgrey} This note summarises the development without
restating the evidence; the full derivations, calibration record and
validation results are held in the firm's internal research documentation.
Figures are in-sample unless stated otherwise; this is not a forecast or an
offer.}

\end{document}
